\section{Гексагональный срез $(2{+}1)$}
\label{sec:hex-slice}

\subsection{Регулярные замощения плоскости}

Локальный закон на плоскости требует замощения без разрывов.
На плоскости существуют ровно три регулярных типа: треугольник, квадрат и шестиугольник
(\cref{fig:carrier-tiling-triangle,fig:carrier-tiling-square,fig:carrier-tiling-hex}).
Число соседей $|N(x)|$ на одном такте --- отдельный вопрос: узлы стоят в центрах плиток
или в вершинах, и квадратная решётка допускает диагональ $\caStepSpace\sqrt{2}$, недостижимую за один~$\htick$
(\cref{fig:carrier-neighbors-tri,fig:carrier-neighbors-sq,fig:carrier-neighbors-hex}).
Наиплотнейшая изотропная упаковка центров даёт гексагональную ячейку Вороного, $|N|=6$
(\cref{fig:carrier-hex-neighbors}).

\begin{figure}[htbp]
  \centering
  \begin{tikzpicture}[carrier panel]
  \foreach \i in {-1,...,1} {
    \foreach \j in {-1,...,1} {
      \pgfmathsetmacro{\x}{\i*0.55 + mod(\j,2)*0.275}
      \pgfmathsetmacro{\y}{\j*0.48}
      \draw[carrier object, carrier fill]
        (\x,\y) -- ++(1,0) -- ++(-0.5,0.866) -- cycle;
    }
  }
  \draw[carrier object, carrier hifill] (0,0) -- (1,0) -- (0.5,0.866) -- cycle;
\end{tikzpicture}

  \caption{Треугольное замощение (выделена одна плитка).}
  \label{fig:carrier-tiling-triangle}
\end{figure}

\begin{figure}[htbp]
  \centering
  \begin{tikzpicture}[carrier panel]
  \foreach \i in {0,...,2} {
    \foreach \j in {0,...,2} {
      \draw[carrier object, carrier fill] (\i,\j) rectangle ++(1,1);
    }
  }
  \draw[carrier object, carrier hifill] (1,1) rectangle ++(1,1);
  \CarrierDimH{1}{2}{-0.35}{$a$}
\end{tikzpicture}

  \caption{Квадратное замощение (выделена одна плитка).}
  \label{fig:carrier-tiling-square}
\end{figure}

\begin{figure}[htbp]
  \centering
  \begin{tikzpicture}[carrier panel]
  \foreach \r in {-1,...,1} {
    \foreach \c in {-1,...,1} {
      \pgfmathsetmacro{\cx}{\c*0.75 + mod(\r,2)*0.375}
      \pgfmathsetmacro{\cy}{\r*0.65}
      \CarrierHex{0.32}{(\cx,\cy)}
    }
  }
  \draw[carrier object, carrier hifill]
    (30:0.32) \foreach \a in {90,150,...,330} { -- (\a:0.32) } -- cycle;
\end{tikzpicture}

  \caption{Шестиугольное замощение (выделена одна плитка).}
  \label{fig:carrier-tiling-hex}
\end{figure}

\begin{figure}[htbp]
  \centering
  \begin{tikzpicture}[carrier panel]
  \foreach \i in {-1,...,1} {
    \foreach \j in {-1,...,1} {
      \pgfmathsetmacro{\x}{\i*0.55 + mod(\j,2)*0.275}
      \pgfmathsetmacro{\y}{\j*0.48}
      \draw[carrier object, carrier fill]
        (\x,\y) -- ++(1,0) -- ++(-0.5,0.866) -- cycle;
    }
  }
  \fill[carrier hifill] (0.5,0.29) circle (2pt);
  \node[carrier site] at (0.5,0.29) {};
\end{tikzpicture}

  \caption{Треугольная решётка: $|N(x)|=3$.}
  \label{fig:carrier-neighbors-tri}
\end{figure}

\begin{figure}[htbp]
  \centering
  \begin{tikzpicture}[carrier panel]
  \foreach \i in {0,...,2} {
    \foreach \j in {0,...,2} {
      \draw[carrier object, carrier fill] (\i,\j) rectangle ++(1,1);
    }
  }
  \foreach \i/\j in {1/1,2/1,0/1,1/2,1/0} {
    \draw[carrier object, carrier hifill] (\i,\j) rectangle ++(1,1);
  }
  \node[carrier site] at (1.5,1.5) {};
  \draw[carrier muted, dashed] (1.5,1.5) -- (2.71,2.71);
  \node[carrier muted, font=\scriptsize] at (2.35,2.45) {$\caStepSpace\sqrt{2}$};
\end{tikzpicture}

  \caption{Квадратная решётка: $|N(x)|=4$; диагональ $\caStepSpace\sqrt{2}$ недостижима за один такт.}
  \label{fig:carrier-neighbors-sq}
\end{figure}

\begin{figure}[htbp]
  \centering
  \begin{tikzpicture}[carrier panel]
  \CarrierHexNeighborhood{0.36}{0.62}
  \node[carrier site] at (0,0) {};
\end{tikzpicture}

  \caption{Гексагональная решётка: $|N(x)|=6$.}
  \label{fig:carrier-neighbors-hex}
\end{figure}

\begin{figure}[htbp]
  \centering
  \begin{tikzpicture}[carrier panel, scale=1.4]
  \CarrierHexRing{0.4}{0.68}
  \filldraw[carrier object] (0,0) circle (2.5pt);
  \draw[carrier object, carrier muted] (0,0) circle (0.4);
  \draw[carrier arrow] (0,0) -- (0.4,0);
  \node[anchor=south, inner sep=1pt] at (0.2,-0.06) {$a=\caStepSpace$};
  \node[anchor=west, inner sep=2pt] at (0.06,0) {$x$};
  \node[anchor=south, inner sep=2pt] at (0,0.92) {$|N(x)|=6$};
\end{tikzpicture}

  \caption{Гексагональная ячейка Вороного: ребро $a=\caStepSpace$ и $|N(x)|=6$.}
  \label{fig:carrier-hex-neighbors}
\end{figure}

\subsection{Константы на гексагональном слое}

Центр правильного шестиугольника и его вершина разделены ребром~$a=\caStepSpace$.
Апофема --- расстояние от центра до стороны --- есть высота равностороннего треугольника со стороной~$a$:
\[
  R_{\mathrm{in}}=\frac{\sqrt{3}}{2}\,a.
\]
Это чистая геометрия однотактового тела, без соседей
(\cref{fig:carrier-hex-radii}).
Отношение к радиусу до вершины $R_{\mathrm{out}}=a$ даёт геометрический множитель среза
\[
  \caKappaGeom_{\mathrm{hex}}=\frac{R_{\mathrm{in}}}{R_{\mathrm{out}}}=\frac{\sqrt{3}}{2}.
\]

\begin{figure}[htbp]
  \centering
  \begin{tikzpicture}[carrier panel, scale=1.5]
  \def\a{1}
  \pgfmathsetmacro{\rin}{0.866*\a}
  \draw[carrier object] (0,0) circle (\a);
  \draw[carrier muted, dashed] (0,0) circle (\rin);
  \draw[carrier object, carrier fill] (0:\a) \foreach \b in {60,120,...,300} { -- (\b:\a) } -- cycle;
  \node[carrier site] at (0,0) {};
  \draw[carrier arrow] (0,0) -- (\a,0);
  \node[anchor=south, inner sep=1pt] at (0.5*\a,-0.06) {$R_{\mathrm{out}}=a$};
  \draw[carrier arrow, carrier muted] (0,0) -- (0,\rin);
  \node[anchor=west, carrier muted, inner sep=1pt] at (0.08,0.5*\rin) {$R_{\mathrm{in}}$};
  \node[anchor=north, font=\footnotesize] at (0,-1.35)
    {$\kappa_{\mathrm{hex}}=R_{\mathrm{in}}/R_{\mathrm{out}}=\sqrt{3}/2$};
\end{tikzpicture}

  \caption{Вписанный и описанный радиусы правильного шестиугольника: $\caKappaGeom_{\mathrm{hex}}$.}
  \label{fig:carrier-hex-radii}
\end{figure}

Тактовая скорость $\caSignalSpeed=\caStepSpace/\htick$ и макроскопическая~$\speedC$ связаны этим же отношением, $\speedC=\caKappaGeom_{\mathrm{hex}} \caSignalSpeed$,
поэтому
\begin{align}
  \htick &= \caKappaGeom_{\mathrm{hex}}\,t_P, \\
  \caSignalSpeed &= \frac{2}{\sqrt{3}}\,c, \\
  \kappa_{\mathrm{link}} &= \frac{1}{|N|} = \frac{1}{6}.
  \label{eq:hex-bridge-constants}
\end{align}
Число~$\speedC$ одно и на срезе, и на трёхмерном носителе, который будет выбран ниже;
меняется только разбиение $\speedC=\caKappaGeom \caSignalSpeed$.

